\backhead{FURTHER EXERCISES}{khred}
\begin{exlist}
\item A sandy loam has pores of equivalent radius 30 $\mu$m and 3 $\mu$m. (a) Find the suction at which each empties ($\gamma=0.0728$ N\,m$^{-1}$, $\beta=0$). (b) Express both in pF. (c) Which of the two still holds water at field capacity (100 cm)?\exans{(a) 49 cm, 495 cm; (b) pF 1.69, 2.69; (c) the 3 $\mu$m pore}
\item Water flows through a single straight tube of radius 50 $\mu$m and length 1 m under a head difference of 1 m ($\mu=1.002\times10^{-3}$ Pa\,s). (a) Find the discharge by Hagen--Poiseuille. (b) By what factor does it change if the radius is halved? (c) How many tubes of 25 $\mu$m carry the same flow as one of 50 $\mu$m?\exans{(a) $2.40\times10^{-11}$ m$^3$\,s$^{-1}$; (b) 1/16; (c) 16}
\item A packing of uniform sand grains of diameter 0.2 mm has porosity 0.40. (a) Find the intrinsic permeability by Kozeny--Carman with $c=180$, $k=d^2\phi^3/[180(1-\phi)^2]$. (b) Convert it to a saturated conductivity in cm\,d$^{-1}$. (c) By what factor does $k$ fall if compaction reduces $\phi$ to 0.35?\exans{(a) $3.95\times10^{-11}$ m$^2$; (b) 3342 cm\,d$^{-1}$; (c) 0.57}
\item A clay soil cracks into planar cracks of aperture 0.2 mm, one crack per 10 cm of horizontal distance. Using the cubic law, find the flux per unit area through the cracks under a unit gradient (cm\,d$^{-1}$), and compare it with the matrix conductivity of 1.7 cm\,d$^{-1}$.\exans{1 cm\,d$^{-1}$; about 0 times the matrix}
\item A loam of porosity 0.50 holds $\theta=0.20$. (a) Find the relative diffusivity of gases by Millington and Quirk. (b) After rain $\theta=0.40$: find it again. (c) By what factor has the diffusion of oxygen fallen?\exans{(a) 0.0723; (b) 0.00186; (c) 39}
\end{exlist}
\backhead{ESSAY QUESTIONS}{khred}
{\small\itshape Answer each question in about one page, with equations, a sketch and at least one number from this chapter.}\par\vspace{3pt}
\begin{exlist}
\item Explain, with the Young--Laplace law, why the retention curve of a soil can be read as a distribution of pore sizes. What assumptions does this reading make, and where does it fail?
\item Discuss the role of throats and bodies in the hysteresis of soil water. Why is the drying curve above the wetting curve, and what does this mean for the reading of a tensiometer after rain?
\item Compare the bundle-of-capillaries model with a network model of the pores. Which properties of a real soil can the network explain that the bundle cannot?
\item Describe how tillage, compaction, roots and earthworms change the geometry of the pore space, and how each change appears in the retention and the conductivity curves.
\item Explain the two stages of evaporation from a bare soil in terms of the pore sizes. Why does a mulch save most water in soils with a long characteristic length?
\item Why does the conductivity of a soil fall by many orders of magnitude while its water content falls by only a factor of two or three? Base your answer on the fourth-power law of flow in a pore.
\end{exlist}
\dykband{EXERCISES IN SWAFL}
{\small Open SWAFL (desktop or online) with the benchmark project. The chapter of the \href{https://github.com/halilunaziru73-creator/Coupled-Water-Fluxes-Pore-Scale-Dynamics-Precision-Irrigation-Mulching-Feedback-Loops/raw/main/Software/v1.1.0/SWAFL_Desktop/docs/manual_v1.1.0/SWAFL_Manual_v1.1.0.pdf}{SWAFL manual} that explains each task is given in brackets; record your values in a table and compare them with what is expected (\emph{Check}).}\par\vspace{3pt}
\begin{exlist}
\item {\color{khblue}[Manual, Chapter 8]} Open the layer \emph{Soil structure and tillage} and choose the sand and then the silty clay loam. For each, read the saturated water content and sketch the retention curve that the layer draws. Which soil has the wider range of pore sizes?\par\noindent{\small\itshape Check: the silty clay loam (smaller $n$), whose curve falls gradually over several decades of suction.}
\item {\color{khblue}[Manual, Chapter 8]} In the same layer move the \emph{Structural consolidation rate} to its smallest and largest value and run. Record how the drainage and the evaporation of the year change, and explain the change with the pore picture of this chapter.\par\noindent{\small\itshape Check: faster consolidation removes the large pores made by tillage, so the drainage of the loosened layer falls.}
\item {\color{khblue}[Manual, Chapter 9]} Open \emph{Mulch and evaporation} and watch the pore picture and the drying front during \emph{Play}. Note the day on which stage 2 begins on the bare soil, with and without mulch.\par\noindent{\small\itshape Check: stage 2 begins later under the mulch; the drying front stays near the surface.}
\item {\color{khblue}[Manual, Chapter 7]} In \emph{Soil, plant and air} compare the water potential of the soil, the root, the leaf and the air on a hot day. Which step has the largest drop, and why?\par\noindent{\small\itshape Check: the step from the leaf to the air, which is far larger than all the others.}
\end{exlist}
